\section{Tier 2: Controlled validation of the filter in SUMO}
\label{sec:phase4-sumo}

\subsection{Motivation}

This tier is the controlled, end-to-end validation of the
partial-observation routing-influence filter in Simulation of Urban
MObility (SUMO) \cite{lopez2018}, the only setting here that supplies
ground-truth regime labels: an
identical population of vehicles is simulated once with sat-navigation
rerouting turned off and once with it turned on. Every component of the
filter is exercised at city-block scale ($n = 2{,}405$ edges): the
inverse-MC estimator is fitted per regime, the log-likelihood-ratio
scorer is applied to held-out trajectories, the feature-regularised
filter variant is tested against a no-features baseline, and the
filter's diagnostic suite (per-row Pearson/cosine, empirical $2$-step
memory audit and the appendix multistart check) is run on the resulting fits.
A trajectory-length sweep traces how the scoring window $T$ shifts both
the empirical reference and the fitted detector together. The heavier
multistart and rush/off-peak proxy audits are summarised in the main text
and reported in Sections~\ref{sec:phase4-multistart} and
\ref{sec:phase4-rho-demand} in the Appendix. The section closes with the
documented difficulty of using the hitting-rate ($f+g$) extension at this
scale, a design-relevant failure mode for the filter's optional $\gobs$
input.

A note on terminology used throughout this section. The no-rerouting
regime is called ``static unguided routing'': SUMO routes those
vehicles by static shortest-path Dijkstra with no live rerouting. This is
a controlled proxy for intrinsic route choice, not a behavioural model of
human preference, and the fitted ``intrinsic'' chain is interpreted
accordingly. The term ``empirical $1$-step reference'' (rather than an
information-theoretic ceiling) means the AUC obtained from the
full-observation transition matrix after Laplace smoothing. Here, Laplace
smoothing means adding a small pseudo-count to each legal transition
before normalising, so legal but unseen transitions do not receive zero
probability. The reference is therefore the best a $1$-step detector can
do \emph{given this finite data and smoothing}, and finite-sample and
train/test variation can move it (and occasionally produce apparent
anomalies, \cref{sec:phase5-multistart}).

\subsection{Method}

\subsubsection{Scenario construction}\label{sec:phase4-scenario}

A rectangular Ingolstadt, Germany study area
($\sim 3.7\,\text{km} \times 3.3\,\text{km}$), derived from
OpenStreetMap (OSM) via \code{netconvert}, yields $n = 2{,}405$ directed edges after filtering
to the largest strongly-connected passenger-allowed component, with
$E = 5{,}987$ directed transitions and mean out-degree $\sim 2.5$. A
demand of $\sim 28{,}800$ vehicles is generated by
\code{randomTrips.py} over a four-hour horizon with \code{-p 0.5}. In
SUMO's \code{randomTrips.py}, \code{-p} is the average departure period in
seconds, so \code{-p 0.5} attempts roughly one new trip every
$0.5\,\mathrm{s}$. A random seed is the integer that initialises this
pseudo-random trip draw; changing it creates a different but comparable
set of origins, destinations and insertion times. Each vehicle is replayed twice
under the same insertion times and origin/destination pairs:
\begin{itemize}
\item \textbf{Intrinsic regime:} \code{--device.rerouting.probability 0}.
Vehicles route via static shortest-path Dijkstra and never react to
congestion.
\item \textbf{Sat-navigation regime:} \code{--device.rerouting.probability 1}.
Every vehicle re-evaluates its route every $30\,\mathrm{s}$ on the basis
of observed live edge travel times.
\end{itemize}
Operating at \code{-p 0.5} was selected from preliminary demand sweeps:
at higher demand both regimes saturate and
the sat-navigation contrast collapses; at lower demand there is
insufficient congestion to drive any rerouting. At \code{-p 0.5} the
satnav regime delivers $\sim 36\%$ more completed trips than intrinsic
with matched insertion ($28{,}098$ vs.\ $20{,}700$ finishing). This
completion imbalance is a limitation of the SUMO operating point. The
validation scorer uses balanced test sets ($N=300$ per regime), and the
Morimura $L_d$ term is scale-free in the aggregate counts, so the result is
not a trivial class-prior effect. Nevertheless, the fitted counts and the
pool of long completed trajectories are conditional on different
completion/survival processes in the two regimes. A real risk is therefore
that the inverse-MC fit learns where trips survive to completion under
sat-navigation rather than only learning routing preferences conditional on
comparable trips. Tier~2 is therefore treated as a controlled operational
validation of the filter under this biased operating point, not as a
completion-balanced causal estimate of the effect of sat-navigation
rerouting. A stricter sensitivity study would rerun the simulation or
rebuild $\fobs$ under matched completed-trip populations; that extension is
left outside the present experimental budget.

For each regime, the experiment exposes vehicle counts \code{edgedata.*.xml} (the
empirical $\fobs$) and complete per-vehicle routes \code{vehroutes.*.xml}
(used as the test set and to construct the empirical chain for the
contrast diagnostic). The inverter sees only $|X_o| = 601$ randomly
	sampled edges, $25\%$ of the network, consistent with the
operating point of the Tier~1 reproduction.

A stationarity caveat applies here and carries through to Tier~3. The
Morimura observation model assumes the counts $\fobs$ are samples of a
\emph{stationary} distribution and that $\gobs$ are stationary hitting
rates. SUMO instead produces finite-horizon, dynamic, congestion-affected
counts over a four-hour rush window. As demand loads, queues form and
rerouting reacts, the instantaneous occupancy distribution changes; the
\code{edgedata} totals are therefore time-integrated visit-frequency
summaries, not exact samples from a single invariant $\pi$. The inverse
fit is consequently a controlled operational test of a stationary surrogate
under structural misspecification, not a claim that the fitted chain is
behavioural ground truth. This is one reason the
full-observation empirical transition matrix is reported as a finite-data reference
rather than as a true ceiling, and why the per-row contrast diagnostic
(\cref{sec:detector}), which compares fitted and empirical \emph{contrasts}
on the same finite data, is the load-bearing evidence rather than the
absolute AUC.

\subsection{Results}

\subsubsection{Single-seed result: features and the inverse fit}
\label{sec:phase4-features-result}

The comparison uses two feature settings. \code{--features none} is the
local-only model, with one parameter per directed transition
($d = n+E = 8{,}392$). \code{--features real} adds shared road-network
features (road class, speed, lane count and turn geometry), giving
$d = 8{,}402$. The \texttt{[all branching]} diagnostic used below means
all unobserved states with more than one legal successor; these are the
rows where the fitted model must generalise beyond directly observed
station counts. On random-trip seed~$42$, $T=20$, $|X_o|=601$ and
$\code{maxiter}=1500$, the headline numbers are summarised in
\cref{tab:phase4-features}.

\begin{figure}[H]
\centering
\includegraphics[width=0.95\linewidth]{fig03_phase4_features_lift.png}
\caption{Tier~2 features lift on the Ingolstadt SUMO network, seed~$42$, $T=20$.
Every metric, including detection AUC, gap-closed, per-row Pearson,
per-row cosine median and fraction of rows with positive cosine, moves
positively under \code{features=real} (dark) vs.\ \code{features=none}
(light). The Pearson +63\% relative lift and the cosine median
$4.6\times$ lift on the \texttt{[all branching]} generalisation block
are the strongest evidence that chain recovery, not aggregate-bias drift,
is what features improve.}
\label{fig:phase4-features-lift}
\end{figure}

\begin{table}[H]
\centering
\caption{Features lift across all metrics on the Ingolstadt SUMO network,
seed=42 (\cref{fig:phase4-features-lift}). \texttt{[all branching]}
restricts the per-row diagnostic to states with
$\mathrm{out\text{-}deg} > 1$ that the inverter did \emph{not} directly
observe, namely the generalisation block.}
\label{tab:phase4-features}
\small
\begin{tabular}{lccc}
\toprule
metric (\texttt{[all branching]}) & \texttt{features=none} & \texttt{features=real} & $\Delta$ \\
\midrule
$\AUC_{\mathrm{emp}}$ (full-obs.\ reference) & $0.685$ & $0.685$ & -- \\
$\AUC_{\mathrm{fit}}$ & $0.612$ & $\mathbf{0.644}$ & $+0.032$ \\
gap closed & $60.6\%$ & $\mathbf{77.5\%}$ & $+16.9$ pp \\
stacked Pearson($\Delta_{\mathrm{fit}}, \Delta_{\mathrm{emp}}$) & $+0.132$ & $\mathbf{+0.215}$ & $+63\%$ rel \\
per-row cosine median & $+0.095$ & $\mathbf{+0.435}$ & $4.6\times$ \\
rows with cosine $> 0$ & $53.5\%$ & $55.4\%$ & $+1.9$ pp \\
\bottomrule
\end{tabular}
\end{table}

Two results are worth separating. First, the $+0.032$ AUC lift on this
seed sits outside the single-instance variability measured in repeated
SUMO baseline runs. As the multi-seed
analysis of \cref{sec:phase4-multiseed} shows, however, this aggregate
AUC lift does \emph{not} survive replication across SUMO instances. The
per-row chain-recovery improvement, not the aggregate AUC, is the robust
effect. Second, the stacked Pearson on \texttt{[all branching]}
rises from $+0.132$ to $+0.215$, and the per-row cosine median from
$+0.095$ to $+0.435$. These are the metrics
that distinguish ``the fit improved because the chains got more accurate
per-row'' from ``the fit improved because the chains drifted in helpful
opposite directions in aggregate''. Both moved positively, supporting the
stronger chain-recovery claim. Unlike the aggregate AUC, this per-row
lift replicates across seeds (\cref{sec:phase4-multiseed}).

The one metric that barely moves is the fraction of rows with positive
cosine (only $+1.9$ pp on \texttt{[all branching]}, $58.6\%$ at
\texttt{[X\_o branching]}). Features sharpen the contrast magnitude in
the rows that were already approximately right, but do not flip the
sign on the $\sim 45\%$ of branching states whose fitted contrast points
the wrong way. This suggests the residual misspecification is
structural to the $1$-step softmax family. \Cref{sec:phase4-2step}
tests that hypothesis.

\subsubsection{Multi-seed validation: chain recovery is robust, detection AUC is not}
\label{sec:phase4-multiseed}

To test whether the seed-$42$ lift replicates, the contrast-correlation
diagnostic was run across five independent SUMO instances
(\code{randomTrips} seeds $\in \{7, 23, 42, 101, 2024\}$, observation-set
seed fixed at $13$, $T=20$, $|X_o|=601$, $\code{maxiter}=1500$) under both
\code{features=none} and \code{features=real}
(\cref{tab:phase4-multiseed}, \cref{fig:phase4-multiseed}).
Because the inverse-MC loss is non-convex, the optimiser can settle into
different local optima, or \emph{basins}, on different SUMO instances. A
basin here means a stable region of parameter space that L-BFGS-B reaches
and then cannot improve beyond, even though another instance may reach a
better region. The multistart audit in the Appendix tests whether random
re-initialisation can escape the low basin.

\begin{figure}[H]
\centering
\includegraphics[width=0.98\linewidth]{fig05_phase4_multiseed.png}
\caption{Multi-seed Tier~2 features sweep (five SUMO seeds,
\texttt{[all branching]} block, $T=20$). \textbf{(a)} The aggregate
detection metric (gap-closed) lifts under \code{features=real} on four
seeds but \emph{reverses} on seed~$101$; the \code{features=none} bars
also reproduce the bimodal basin structure ($\{7,23\}$ low,
$\{42,101,2024\}$ high). \textbf{(b)} The per-row chain-recovery metric
(stacked Pearson) lifts on \emph{all five} seeds. Features are a robust
chain-recovery lever, not a robust aggregate-detection lever.}
\label{fig:phase4-multiseed}
\end{figure}

\begin{table}[H]
\centering
\caption{Multi-seed features sweep on the Ingolstadt SUMO network
(\cref{fig:phase4-multiseed}), \texttt{[all branching]} block,
$T=20$, $|X_o|=601$. The AUC lift reverses on seed~$101$; the Pearson
lift is positive on all five seeds. Mean~$\Delta$ rows are
\code{real}~$-$~\code{none} with sample standard deviation.}
\label{tab:phase4-multiseed}
\small
\begin{tabular}{lcccccc}
\toprule
 & \multicolumn{2}{c}{gap closed} & \multicolumn{2}{c}{$\AUC_{\mathrm{fit}}$} & \multicolumn{2}{c}{Pearson} \\
\cmidrule(lr){2-3}\cmidrule(lr){4-5}\cmidrule(lr){6-7}
seed & none & real & none & real & none & real \\
\midrule
$7$    & $30.3\%$ & $60.3\%$          & $0.544$ & $0.588$          & $+0.191$ & $+0.238$ \\
$23$   & $34.2\%$ & $53.8\%$          & $0.554$ & $0.585$          & $+0.133$ & $+0.194$ \\
$42$   & $60.6\%$ & $77.5\%$          & $0.612$ & $0.644$          & $+0.132$ & $+0.215$ \\
$101$  & $62.1\%$ & $\mathbf{45.0\%}$ & $0.581$ & $\mathbf{0.559}$ & $+0.171$ & $+0.194$ \\
$2024$ & $61.8\%$ & $85.1\%$          & $0.585$ & $0.618$          & $+0.131$ & $+0.168$ \\
\midrule
mean $\Delta$ & \multicolumn{2}{c}{--} & \multicolumn{2}{c}{$+0.024 \pm 0.026$} & \multicolumn{2}{c}{$\mathbf{+0.050 \pm 0.023}$} \\
\bottomrule
\end{tabular}
\end{table}

Three findings. First, the aggregate detection lift is \emph{not}
statistically robust: the paired $\AUC_{\mathrm{fit}}$ lift averages
$+0.024 \pm 0.026$ across the five seeds, positive on four but reversing
on seed~$101$ (a paired $t$-test gives $t \approx 2.0$, $p \approx 0.11$,
not significant at $\alpha = 0.05$). The single-seed $+0.032$ headline of
\cref{sec:phase4-features-result} is therefore a favourable draw, not a
replicable property of the aggregate AUC.

Second, the per-row chain-recovery metrics \emph{are} robust. The stacked
Pearson lift averages $+0.050 \pm 0.023$ and is positive on \emph{all five} seeds
($t \approx 4.9$, $p \approx 0.008$); the per-row cosine median lifts on
all five as well (mean $\approx +0.25$). Seed~$101$ is the clearest warning
that AUC and chain quality can move in different directions: features
\emph{hurt} its detection AUC ($0.581 \to 0.559$) while \emph{improving}
its chain recovery (cosine median $0.062 \to 0.475$). The features lever
is therefore a
\textbf{chain-recovery} improvement, not an aggregate-detection
improvement.

Third, the \code{features=none} fit reproduces the bimodal gap-closure
structure observed in repeated SUMO runs: seeds $\{7, 23\}$ cluster
low ($30$--$34\%$) and seeds $\{42, 101, 2024\}$ cluster high
($61$--$62\%$), with nothing in between. The same seeds fall in the same
clusters as the earlier run, indicating that the low/high split is a
property of the loss landscapes on these instances rather than measurement
noise. The \code{features=none} $\AUC_{\mathrm{fit}}$ across seeds is
$0.575 \pm 0.027$, consistent with the baseline variability. A follow-up
multistart audit tested whether the low basin could be escaped by random
re-initialisation; the detailed audit is reported in
Section~\ref{sec:phase4-multistart} in the Appendix, and it did not escape
the low basin.

\subsubsection{Higher-order memory diagnostic: a regularised 2-step
empirical lifts the reference}
\label{sec:phase4-2step}

A natural hypothesis for the still-flipped rows is that the $1$-step
Markov assumption is too forgetful. The fitted model sees only the current
edge $x_t$, but drivers may also condition on the edge they just came from:
``do not U-turn'', ``continue along the arterial'', or ``stay in this
lane group''. A $2$-step model gives the next-edge distribution the context
$(x_{t-1},x_t)$ instead of only $x_t$.

This subsection is a diagnostic, not a new fitted inverse-MC model. The
test is: if a full-observation empirical $2$-step scorer ranks the two
SUMO regimes better than the full-observation empirical $1$-step scorer,
then the data contain route-memory signal that the current $1$-step
softmax family cannot express. If it does not lift, the flipped rows are
more likely caused by optimisation, sparse observation, or feature
misspecification.

The catch is sparsity. On the Ingolstadt network the raw context space
$(x_{t-1},x_t)$ has roughly $|X|^2 \approx 5.8$M possible pairs, and most
possible pairs are never seen. A raw $2$-step empirical matrix would
therefore overfit rare contexts. Two regularisation schemes are tested
across the same five SUMO seeds used in \cref{sec:phase4-multiseed}
(training/test split unchanged):

\paragraph{Dirichlet smoothing: shrink rare contexts toward uniform turns.}
\[
\widehat{P}_{2,\alpha}(y | x_p, x) =
\frac{c(x_p, x, y) + \alpha}
{c(x_p, x) + \alpha \cdot |\mathrm{adj}(x)|},
\]
swept over $\alpha \in \{10^{-3}, 10^{-2}, 10^{-1}, 1, 10, 100\}$.

\paragraph{Interpolation backoff: shrink rare contexts toward the 1-step empirical.}
\[
\widehat{P}_{\mathrm{back}}(y | x_p, x) =
\lambda \cdot \widehat{P}_{2}^{\mathrm{MLE}}(y | x_p, x)
+ (1 - \lambda) \cdot \widehat{P}_{1}(y | x),
\]
where $\widehat{P}_{2}^{\mathrm{MLE}}$ is the unsmoothed 2-step MLE
and $\widehat{P}_{1}$ is the 1-step empirical chain that elsewhere
serves as the full-observation reference. Two $\lambda$ modes are tested:
\begin{itemize}
\item fixed $\lambda \in \{0, 0.25, 0.5, 0.75, 0.9, 1.0\}$;
\item count-conditional
$\lambda = N(x_p, x) / (N(x_p, x) + \kappa)$ with
$\kappa \in \{1, 10, 100, 1000\}$.
\end{itemize}
The count-conditional mode is the intuitive version: trust the $2$-step
estimate when a context has been seen often, and fall back smoothly to the
$1$-step empirical for rare contexts.

\begin{table}[H]
\centering
\caption{2-step empirical sweep on the Ingolstadt SUMO network across five SUMO seeds
($T=20$, training/test split as in \cref{sec:phase4-multiseed}, 1-step
empirical baseline mean $0.652 \pm 0.022$ across these seeds).
Each row reports the best configuration for that method. The single-seed
row is a comparison point; the multi-seed rows are the load-bearing
evidence.}
\label{tab:phase4-2step}
\scriptsize
\begin{tabular}{@{}p{0.27\linewidth}p{0.18\linewidth}ccc@{}}
\toprule
method & config & $\AUC_{\mathrm{2step}}$ & $\Delta$ & wins \\
\midrule
single-seed Dirichlet, seed~42 & $\alpha \to 0$ & $0.666$ & $-0.019$ & $0/1$ \\
\midrule
Dirichlet, multi-seed            & $\alpha = 10^{-3}$ & $0.673 \pm 0.012$ & $+0.021$ & $4/5$ \\
fixed backoff, multi-seed        & $\lambda = 0.9$    & $0.747 \pm 0.012$ & $+0.095$ & $5/5$ \\
\textbf{count backoff, multi-seed} & $\boldsymbol{\kappa = 100}$ & $\boldsymbol{0.752 \pm 0.015}$ & $\boldsymbol{+0.100}$ & $\boldsymbol{5/5}$ \\
\bottomrule
\end{tabular}
\end{table}

Three findings (\cref{tab:phase4-2step}).

First, the single-seed Dirichlet result is misleading. On seed~42 it gives
$0.666$, below the $1$-step reference of $0.685$, but seed~42 is also the
highest-reference seed in the sweep. There is less room to improve on
that seed than on the others. More importantly, Dirichlet smoothing
shrinks rare $2$-step contexts toward uniform outgoing turns, which throws
away useful $1$-step routing information.

Second, backoff to the $1$-step empirical is the right shrinkage target.
At $\kappa = 100$, the count-conditional rule gives positive lifts on all
five seeds, ranging from $+0.090$ to $+0.110$ (mean $+0.100$, standard
deviation $0.009$). On seed~42 the effective $\lambda$ across the
$3{,}528$ observed $(x_p,x)$ contexts has median $0.21$ and 10th--90th
percentiles $[0.03,0.81]$. In plain terms, the scorer mostly uses the
trusted $1$-step baseline and only lets the $2$-step context dominate
where the data support it.

Third, the conclusion is not sensitive to one hand-picked hyperparameter.
The fixed-$\lambda$ sweep gives lifts of
$+0.088$ to $+0.095$ for $\lambda \in \{0.25,0.5,0.75,0.9\}$, all
positive across all five seeds. The count-conditional sweep gives
$+0.095$ to $+0.100$ for $\kappa \in \{1,10,100,1000\}$. The exact
$+0.100$ maximum is mildly test-tuned, but the qualitative conclusion is
stable: any reasonable backoff strength lifts AUC by about $+0.09$ to
$+0.10$ across all five SUMO seeds.

\paragraph{Implication.} A properly regularised 2-step empirical
\emph{does} lift the reference, decisively, across every SUMO seed
tested: the regularised 2-step AUC reaches $0.752 \pm 0.015$ against
the 1-step baseline of $0.652 \pm 0.022$, a $+0.100$ delta whose
per-seed signs all point the same way. Two implications follow.

The per-row chain-recovery gap left by the $1$-step features lever
(\cref{sec:phase4-features-result}: $+0.215$ stacked Pearson on
\texttt{[all branching]}, $\sim 45\%$ of rows still with
wrong-direction contrast) is consistent with the $1$-step softmax
family being structurally unable to express $(x_{t-1}, x_t)$-conditional
choice. This is the regime where a parametric 2-step Morimura
extension is the right design lever, with the regularised empirical
sitting roughly $+0.10$ AUC above the 1-step reference at $T = 20$ as the
performance target it would need to approach.

No $\AUC_{\mathrm{fit}}$ is reported for the regularised
2-step diagnostic. A parametric 2-step Morimura inverter is a non-trivial
extension with a different state space, different gradient and different
reference argument, and is left to future work
(\cref{sec:disc-future}). The empirical result above establishes that
the signal exists and bounds the achievable 2-step inverter
performance from above; what it does not do is settle whether the
parametric inversion route can recover that bound under partial
observation.

\subsubsection{Trajectory-length sweep: lifting the reference and the detector together}
\label{sec:phase4-tsweep}

The $0.685$ reference is specific to the chosen test trajectory length
$T = 20$. The log-likelihood-ratio score $\Lam(\tau)$ is a sum of $T$
independent log-ratio increments, so its variance scales as $\sim T$.
Longer trajectories produce a tighter score distribution between the
two regimes and a higher achievable AUC; the data are ``the same'',
only the observation window is wider.

\begin{figure}[H]
\centering
\includegraphics[width=0.78\linewidth]{fig02_phase4_tsweep.png}
\caption{Tier~2 $T$-sweep on the Ingolstadt SUMO network, seed~$42$,
\code{features=real}. The empirical $1$-step Markov reference
(grey dashed) rises with $T$ as expected for a variance-shrinking
score; the fitted detector (blue solid) tracks it with roughly constant
gap. The shaded region is the recoverable contrast not yet captured by
the fitter. The fitted chain itself is $T$-independent (identical final
losses across runs); only the trajectory-level scoring window changes.}
\label{fig:phase4-tsweep}
\end{figure}

\begin{table}[H]
\centering
\caption{$T$-sweep on the Ingolstadt SUMO network, seed=42, \code{features=real}
(\cref{fig:phase4-tsweep}). The fit itself is $T$-independent: only
the test scoring window changes, as confirmed by the identical
final-loss values $836.4 / 529.4$ across all three runs.}
\label{tab:phase4-tsweep}
\small
\begin{tabular}{ccccc}
\toprule
$T$ & $\AUC_{\mathrm{emp}}$ & $\AUC_{\mathrm{fit}}$ & gap closed & Pearson (all-br) \\
\midrule
$20$ & $0.685$ & $0.644$ & $77.5\%$ & $+0.215$ \\
$30$ & $0.763$ & $0.677$ & $67.4\%$ & $+0.220$ \\
$40$ & $\mathbf{0.801}$ & $\mathbf{0.708}$ & $69.0\%$ & $+0.216$ \\
\bottomrule
\end{tabular}
\end{table}

Two findings worth recording. First, the empirical reference rises from
$0.685$ to $0.801$ between $T=20$ and $T=40$, confirming that
the reference at $T=20$ was an artefact of the truncation window, not a
fundamental data limit. Second, the per-row Pearson is constant across
$T$ at $\sim +0.21$. $T$ is purely a scoring-window lever: it lets the
trajectory-level score discriminate more sharply, but it does not
change the chain that was fit. The fitted chain captures the same
$\sim 21\%$ of the true per-row contrast regardless of how long each test
trip is observed.

The strongest single-seed SUMO headline is therefore $\AUC = 0.708$
at $T=40$, $69\%$ of the gap closed against the full-observation
empirical $1$-step reference of $0.801$, with a stable per-row chain
Pearson of $+0.216$ on \texttt{[all branching]}.

A separate simulator stress test varied sat-navigation adoption and demand
independently before the rush-vs-off-peak proxy was used in Tier~3. The
full $\rho \times d$ decomposition is reported in
Section~\ref{sec:phase4-rho-demand} in the Appendix. In brief, under
SUMO's static unaided-driver baseline, changing demand alone produced a
near-random empirical AUC ($0.493 \pm 0.018$), while adding full
sat-navigation at peak demand produced a large signal
($0.780 \pm 0.013$). This supports the rush/off-peak proxy only as a weak
simulator audit: real commuters can remember recurring congestion, so the
real-world proxy may still mix human experience with algorithmic guidance.

\subsubsection{The \texorpdfstring{$f+g$}{f+g} extension at scale: a refined diagnosis}
\label{sec:phase4-fg-infeasible}

The optional $f+g$ loss was audited because the initial SUMO-scale run at
$|X_o|=601$ did not terminate after $9.5$ hours, but the diagnosis is more
specific than simple infeasibility. A finite-difference-verified sparse-LU
implementation confirms the mathematics, yet naive sparse solves are still
slower than dense at full SUMO scale because the wide gradient solve
dominates. On smaller Xuancheng subgraphs the solver is fast enough, but
a different problem appears: most observed station pairs are never hit, so
their empirical $g$ values sit at the same small floor. Intermediate
$\gamma$ then lets this low-information $g$ term overwhelm the useful
count signal $f$; only high-$\gamma$ settings close to the $f$-only end
recover the empirical reference. The practical detector in this
dissertation therefore remains $f$-only, and the solver profiling plus
$\gamma$ sweep are moved to Section~\ref{app:fg-profiling} in the
Appendix.

\subsection{Interpretation}

Tier~2 is the positive validation tier, but only under the controlled
conditions that make its labels meaningful. With ground-truth rerouting
labels and matched OD pairs, the fitted filter discriminates sat-navigation
from static unguided routing, reaching $\AUC=0.708$ at $T=40$ while
closing $69\%$ of the gap to the full-observation empirical $1$-step
reference. The feature-augmented score is useful primarily as a
chain-recovery lever: it improves per-row Pearson/cosine consistently
across five seeds, but its aggregate AUC lift is not robust.

The tier also identifies the model-class boundary. A regularised
empirical $2$-step reference lifts AUC by about $+0.09$--$+0.10$ on every
SUMO seed, showing that higher-order route memory is a real signal in the
controlled data and that a parametric $2$-step inverse-MC model is the
next justified detector design. By contrast, the $f+g$ extension is not
ready at city scale under the current solver: subgraph tests show that it
can help when $\gamma$ is tuned near the f-only end, but full-network
runtime and hitting-rate sparsity keep the deployed Tier~2 detector
f-only. The defensible claim from this tier is therefore specific:
the current $1$-step partial-observation filter works as a controlled
labelled detector, and its diagnostics correctly expose where richer
model structure is needed.
